% Pregunta 884 (2003-julio-OEP2003, n.º 12): maximin de Rawls frente a Pareto.
% E1 = (U_B, U_A) = (1, 10) y E2 = (1, 20) están en la misma curva en L de W = min{U_A, U_B} = 1
% (vértice en (1,1)): el maximin es indiferente, aunque E2 está en el cuadrante de
% mejoras de Pareto de E1. A escala: 1 unidad de U_B = 1,1 cm; 1 unidad de U_A = 0,22 cm.
% Autor: Víctor Gutiérrez Marcos
\documentclass[tikz,border=6pt]{standalone}
\input{tikz-tcee}
\begin{document}
\begin{tikzpicture}[tcee,x=1.1cm,y=0.22cm]
  % Cuadrante de mejoras de Pareto desde E1 (U_B >= 1, U_A >= 10)
  \fill[area] (1,10) rectangle (5.4,24);
  \node[etiqueta,text=tceemorado,align=center] at (3.6,15) {mejoras de Pareto\\desde $E_1$};
  \draw[eje] (0,0) -- (5.8,0) node[right] {$U_B$};
  \draw[eje] (0,0) -- (0,25.5) node[above] {$U_A$};
  \node[below left] at (0,0) {$0$};
  \foreach \t in {1,...,5} \draw (\t,0) -- (\t,-0.4) node[below] {$\t$};
  \foreach \t in {10,20} \draw (0,\t) -- (-0.08,\t) node[left] {$\t$};
  % Curva de indiferencia social de Rawls que pasa por E1 y E2
  \draw[curva1] (1,24) -- (1,1) -- (5.4,1) node[etiqueta,above left,text=tceemorado] {$W=\min\{U_A,U_B\}=1$};
  % Proyecciones
  \draw[guia] (0,10) -- (1,10);
  \draw[guia] (0,20) -- (1,20);
  % Paso de E1 a E2
  \draw[flecha] (1.25,10.8) -- (1.25,19.2);
  \node[etiqueta,right,text=tceerojo,align=left] at (1.3,22.2) {A gana, B igual:\\Pareto prefiere $E_2$};
  \node[punto] at (1,10) {};
  \node[etiqueta,left] at (0.95,11.6) {$E_1$};
  \node[punto] at (1,20) {};
  \node[etiqueta,left] at (0.95,21.6) {$E_2$};
  \node[etiqueta,text=tceemorado,align=left,right] at (1.3,7) {misma L: Rawls, indiferente};
\end{tikzpicture}
\end{document}
